\newpage
\section{\MakeUppercase{Conclusion}}

The aim of this project was to design and implement a modular bilingual speech-to-dialogue pipeline, which separates the functionality of standard components of a conversational AI system such that each module can be built, optimised, and swapped out independently. The findings achieved across all components confirm this architectural design, and that an intelligently written modular system can provide competitive performance without compromising either flexibility or maintainability.

An English ASR model based on the FastConformer-CTC-BPE architecture, fine-tuned using parameter-efficient adapter modules, attained a WER of 6.14\% on the test set, an impressive result given that only 0.38\% of the total model parameters were updated during fine-tuning. This supports the finding that PEFT-based adaptation of pre-trained Conformer models is a viable and resource-efficient approach to domain-specific ASR. The Nepali ASR, built on the fine-tuned Quantized Indic Conformer, achieved a WER of 31.40\% on the clean FLEURS test set, outperforming all evaluated Whisper variants despite having fewer parameters. Under noisy conditions, the model attained a WER of 33.75\%, demonstrating meaningful noise robustness. The ONNX variant with quantisation further compressed the model to 133.94~MB, making it sufficiently compact for resource-constrained deployment at an acceptable accuracy trade-off.

The small language model, a self-trained decoder-only Transformer, achieved a perplexity of 1.96, a BERTScore-F1 of 0.9303, and a ROUGE-L of 0.5862, indicating that the model is both lightweight and capable of producing semantically coherent and contextually appropriate responses. RAG integration using a FAISS-backed knowledge base of recent news and events significantly improved the factual grounding and relevance of answers generated by the SLM in isolation, directly addressing the hallucination and knowledge staleness inherent to parametric language models trained on fixed datasets.

Together, the bilingual ASR pathways, the RAG-enhanced SLM, and the parallel English and Nepali text-to-speech modules form a complete, end-to-end spoken dialogue system capable of accepting speech in either language and returning contextually relevant, up-to-date spoken responses in the user's preferred language. The system addresses a genuine gap in conversational AI accessibility for Nepali-speaking users and demonstrates that low-resource Indic languages can be effectively integrated into a production-focused pipeline through careful model selection, fine-tuning, and quantisation.

The modular design philosophy adopted in this project represents one of its most significant strengths. Decoupling the components allowed the ASR models to be fine-tuned independently, the knowledge base to be updated without retraining any model, and each stage to be evaluated in isolation. This architecture is conducive to continuous improvement: as stronger alternatives emerge, individual components can be replaced without requiring a full system retraining.

In conclusion, this project demonstrates that a well-structured, decoupled, modular pipeline is not only architecturally sound but also practically viable, delivering strong quantitative performance across ASR, language modelling, and retrieval, while remaining flexible, maintainable, and extensible for future development in real-world multilingual conversational AI applications.
